% SHARD-ID: AQM-63-EVALUATED-DERHAM-CONSTRAINTS-TANGENT-WEYL
% SHARD-TITLE: Evaluated de Rham Constraints on Tangent-Weyl Kinematics
% SHARD-SUMMARY: Reformulates exact de Rham one-forms as cotangent Weyl constraints on \(\ell^2(TX)\).
% SHARD-SUMMARY: Proves the compatibility statement and separates it from a full CSS rule requiring extra tangent labels.
% SHARD-SUMMARY: Works through finite examples showing that unrestricted exact forms often collapse finite rational supports.
% SHARD-KEYWORDS: de Rham, tangent Weyl, cotangent constraints, stabilizer, F3, circle, affine line, crossing

\section{Evaluated de Rham Constraints on Tangent-Weyl Kinematics}
\label{sec:evaluated-derham-constraints-tangent-weyl}

\subsection{Status and Source Anchors}

This shard refines the compatibility question left open by AQM-60--AQM-62.
The tangent-cotangent Hilbert space is the AQM-61 convention
\(\mathcal H_P^{\rm geom}=\ell^2(T_PX)\).  The de Rham source is AQM-54:
exact one-forms come from \(d:A\to\Omega^1_{A/k}\), and \(d_1d_0=0\) gives
the old finite CSS row-orthogonality test.  The present shard keeps only the
part which is canonical for the new kinematics: evaluate exact one-forms at
finite-residue points and use them as cotangent, hence \(Z\)-type, Weyl labels.

\subsection{Definition: Evaluated Exact-One-Form Constraints}

Let \(X=\Spec A\), and let \(S\) be a finite set of finite-residue points.  Put
\[
  K_P=\kappa(P),\qquad V_P=T_PX,\qquad V_P^*=T_P^*X
\]
and
\[
  V_S=\bigoplus_{P\in S}V_P,\qquad
  V_S^*=\bigoplus_{P\in S}V_P^*,\qquad
  \mathcal H_S=\ell^2(V_S).
\]
For \(\omega\in\Omega^1_{A/k}\), define its fibrewise evaluation
\[
  \operatorname{ev}^1_S(\omega)=
  (\omega\otimes1)_P{}_{P\in S}\in V_S^*.
\]
The compatible de Rham momentum label space is
\[
  L_Z^{\rm dR}(S)=\operatorname{ev}^1_S(dA)\subset V_S^*.
\]
For \(\alpha=(\alpha_P)\in V_S^*\), the Weyl phase operator is
\[
  R(\alpha)|v\rangle
  =
  \prod_{P\in S}\psi_{K_P}(\alpha_P(v_P))\,|v\rangle,
  \qquad v=(v_P)\in V_S.
\]
All such \(R(\alpha)\) commute with each other.  Therefore
\[
  \{R(\alpha):\alpha\in L_Z^{\rm dR}(S)\}
\]
is a valid commuting stabilizer family after the \(+1\) phases are chosen.

The joint \(+1\) space is explicit.  Define the annihilator
\[
  A_S^{\rm dR}=
  \{v\in V_S:
    \prod_{P\in S}\psi_{K_P}(\alpha_P(v_P))=1
    \text{ for every }\alpha\in L_Z^{\rm dR}(S)\}.
\]
Then
\[
  \mathcal H_S^{Z{\rm dR}}=\ell^2(A_S^{\rm dR})\subset\ell^2(V_S).
\]
When all points are \(k=\mathbb F_p\)-rational, this is the simpler condition
\[
  \sum_{P\in S}\alpha_P(v_P)=0
  \qquad\text{for all }\alpha\in L_Z^{\rm dR}(S).
\]

\subsection{Why This Is Compatible but Not Yet Full CSS}

The construction is compatible because exact one-forms evaluate to cotangent
vectors, and cotangent vectors are exactly the phase labels in
\[
  T_PX\oplus T_P^*X.
\]
No basis of \(\Omega^1_{A/k}\) is needed, and no artificial copy of
\(\Omega^1\) is used as a tangent configuration space.

It is not yet the old full de Rham CSS rule.  An \(X\)-type label must be a
tangent vector \(w\in V_S\).  It commutes with all de Rham \(Z\)-constraints
exactly when
\[
  w\in A_S^{\rm dR}.
\]
Thus any chosen subspace
\[
  L_X\subset A_S^{\rm dR}
\]
gives a genuine CSS completion, but this \(L_X\) is extra data.  Possible
sources are a metric, Hodge star, action functional, Lagrangian correspondence,
or a deliberately chosen subspace of the annihilator.  Without such data the
canonical output is the \(Z\)-constraint space \(\ell^2(A_S^{\rm dR})\), not a
full dynamical code.

\subsection{Example 1: A Reduced Point}

Let \(A=k=\mathbb F_3\).  Then \(\Omega^1_{k/k}=0\), \(V_P=0\), and
\[
  \mathcal H_P=\ell^2(0)\cong\mathbb C.
\]
Also \(dA=0\), so
\[
  L_Z^{\rm dR}=0,\qquad A_P^{\rm dR}=0.
\]
The evaluated de Rham constraint does nothing, as it should.

\subsection{Example 2: Dual Numbers over F3}

Let
\[
  A=k[\epsilon]/(\epsilon^2),\qquad k=\mathbb F_3,\qquad P=(\epsilon).
\]
AQM-61 gives
\[
  V_P=k\,\partial_\epsilon,\qquad V_P^*=k\,d\epsilon,\qquad
  \mathcal H_P=\ell^2(k)\cong\mathbb C^3.
\]
For \(f=a+b\epsilon\),
\[
  df=b\,d\epsilon.
\]
Therefore
\[
  L_Z^{\rm dR}=k\,d\epsilon=V_P^*.
\]
For \(v=r\partial_\epsilon\),
\[
  d\epsilon(v)=r.
\]
The annihilator condition is \(br=0\) for all \(b\in k\), hence \(r=0\).  Thus
\[
  A_P^{\rm dR}=0,\qquad
  \mathcal H_P^{Z{\rm dR}}=\mathbb C|0\rangle.
\]
This is compatible with tangent-Weyl kinematics, but it is a \(Z\)-phase
constraint.  The old AQM-54 \(X\)-check is its Fourier-dual presentation, not
the canonical tangent-coordinate reading.

\subsection{Example 3: The Rational Affine Line over F3}

Let \(X=\Spec k[x]\), \(k=\mathbb F_3\), and take rational support
\[
  S=\{0,1,2\}.
\]
AQM-62 gives
\[
  V_S=k^3,\qquad \mathcal H_S\cong(\mathbb C^3)^{\otimes3}.
\]
The exact one-form \(df=f'(x)\,dx\) evaluates to
\[
  (f'(0),f'(1),f'(2))\in k^3=V_S^*.
\]
The three functions \(x,x^2,x^5\) give
\[
\begin{array}{c|c}
f & (f'(0),f'(1),f'(2))\\ \hline
x & (1,1,1)\\
x^2 & (0,2,1)\\
x^5 & (0,2,2).
\end{array}
\]
The determinant of these three rows is \(2\in\mathbb F_3^\times\), so
\[
  L_Z^{\rm dR}(S)=V_S^*.
\]
Consequently
\[
  A_S^{\rm dR}=0,\qquad
  \mathcal H_S^{Z{\rm dR}}=\mathbb C|0,0,0\rangle.
\]
This is the first warning: unrestricted exact forms on a finite rational
support can be too strong as a dynamical selection principle.

\subsection{Example 4: Two Independent Dual Directions}

Let
\[
  A=k[u,v]/(u^2,v^2),\qquad k=\mathbb F_3,\qquad P=(u,v).
\]
AQM-61 computes
\[
  V_P=k\,\partial_u\oplus k\,\partial_v,\qquad
  \mathcal H_P=\ell^2(k^2)\cong\mathbb C^9.
\]
The exact forms \(du\) and \(dv\) occur as \(d(u)\) and \(d(v)\).  Hence
\[
  L_Z^{\rm dR}=k\,du\oplus k\,dv=V_P^*.
\]
The annihilator is zero, so
\[
  \mathcal H_P^{Z{\rm dR}}=\mathbb C|0,0\rangle.
\]
Again the construction is compatible, but unrestricted exact one-forms freeze
all tangent coordinates.

\subsection{Example 5: The Crossing over F3}

Let \(A=k[x,y]/(xy)\) and take the rational support.  AQM-62 gives one
two-dimensional tangent at the origin and four one-dimensional tangents away
from it:
\[
  V_S\cong k^6,\qquad
  \mathcal H_S\cong\mathbb C^9\otimes(\mathbb C^3)^{\otimes4}.
\]
A function on the crossing is represented by an \(x\)-axis polynomial and a
\(y\)-axis polynomial with the same constant term.  Its differential along the
\(x\)-axis gives arbitrary derivative values at \(x=0,1,2\), by Example 3.
Its differential along the \(y\)-axis gives arbitrary derivative values at
\(y=0,1,2\), again by Example 3.  These six derivatives are exactly the six
cotangent coordinates dual to \(V_S\).  Hence
\[
  L_Z^{\rm dR}(S)=V_S^*,\qquad A_S^{\rm dR}=0,
\]
and
\[
  \mathcal H_S^{Z{\rm dR}}=\mathbb C|0,0,0,0,0,0\rangle.
\]
The singular tangent at the origin is real kinematics, but exact one-forms on
the full rational support still kill it.

\subsection{Example 6: The Circle over F3}

Let
\[
  C=\Spec k[x,y]/(x^2+y^2-1),\qquad k=\mathbb F_3.
\]
Use the rational point order
\[
  (1,0),(2,0),(0,1),(0,2).
\]
AQM-62 gives
\[
  V_S\cong k^4,\qquad \mathcal H_S\cong(\mathbb C^3)^{\otimes4}.
\]
Use tangent bases \((0,1)\) at the first two points and \((1,0)\) at the last
two.  The following exact forms evaluate to cotangent rows:
\[
\begin{array}{c|c}
f & (df_P(\tau_P))_{P\in S}\\ \hline
y & (1,1,0,0)\\
x & (0,0,1,1)\\
xy & (1,2,1,2)\\
xy^3 & (0,0,1,2).
\end{array}
\]
The determinant of this \(4\times4\) row matrix is \(2\), so the rows span
\[
  V_S^*.
\]
Thus
\[
  L_Z^{\rm dR}(S)=V_S^*,\qquad
  \mathcal H_S^{Z{\rm dR}}=\mathbb C|0,0,0,0\rangle.
\]
So on the finite rational circle, the naive unrestricted exact-form principle
again selects the zero tangent state.

\subsection{Conclusion}

The compatibility result is positive but austere:
\[
  dA\quad\leadsto\quad L_Z^{\rm dR}(S)\subset
  \bigoplus_{P\in S}T_P^*X
\]
is canonical and acts on the tangent-Weyl Hilbert space.  However, the examples
show that using all exact forms on a finite rational support often gives
\[
  L_Z^{\rm dR}(S)=V_S^*
\]
and therefore a one-dimensional zero-tangent state.  A nontrivial dynamical
principle must restrict the allowed exact forms, add an action/Hamiltonian, or
choose additional \(X\)-type tangent constraints from a justified source.  The
old de Rham CSS test is therefore compatible only after this fibrewise
reinterpretation, and it remains a prototype rather than a final dynamics.
